% Part: sets-functions-relations
% Chapter: arithmetization
% Section: cuts
% Hindi translation (hi-Deva-IN) of Open Logic commit 9620cc73f9c8e0ad003c514a5d3748f29611c4c0.
%
\documentclass[../../../include/open-logic-section]{subfiles}

\begin{document}

\olfileid[hi]{sfr}{arith}{cuts}
\olsection{$\Rat$ से $\Real$ तक}

मूलतः पूर्णता गुणधर्म दिखाता है कि वास्तविक संख्या-रेखा का प्रत्येक बिंदु
$\alpha$ उस रेखा को ठीक दो भागों में बाँटता है: एक वह भाग जिसका लघुतम ऊपरी
परिबंध $\alpha$ है, और दूसरा वह भाग जिसका महत्तम निचला परिबंध $\alpha$ है।
परिमेय संख्याओं से वास्तविक संख्याओं का \emph{निर्माण} करने के लिए डेडेकिंड ने
सुझाव दिया कि वास्तविक संख्याओं को बस उन \emph{विच्छेदों} के रूप में मानें जो
परिमेय संख्याओं का विभाजन करते हैं। अर्थात हम $\sqrt{2}$ को उस \emph{विच्छेद}
से पहचानते हैं जो $< \sqrt{2}$ वाली परिमेय संख्याओं को $> \sqrt{2}$ वाली
परिमेय संख्याओं से अलग करता है।

इसे अधिक व्यवस्थित करें। परिमेय संख्याओं को दो भागों में काटने पर, बने हुए
विभाजन के केवल \emph{निचले} भाग को देखकर भी उसे अद्वितीय रूप से पहचाना जा सकता
है। इसलिए सटीक रूप में यह परिभाषा देते हैं:

\begin{defn}[विच्छेद]
कोई \emph{विच्छेद} $\alpha$ परिमेय संख्याओं का ऐसा रिक्तेतर उचित आरंभिक खंड
है जिसका कोई अधिकतम अवयव नहीं होता। अर्थात $\alpha$ विच्छेद तभी और केवल तभी
है जब:
\begin{enumerate}
	\item \emph{रिक्तेतर और उचित}: $\emptyset \neq \alpha \subsetneq \Rat$
	\item \emph{आरंभिक}: सभी $p,q \in \Rat$ के लिए, यदि $p < q \in \alpha$ तो $p \in \alpha$
	\item \emph{कोई अधिकतम नहीं}: प्रत्येक $p \in \alpha$ के लिए कोई $q \in \alpha$ ऐसा है कि $p < q$
\end{enumerate}
तब $\Real$ विच्छेदों का समुच्चय है।
\end{defn}

अब कह सकते हैं कि $\sqrt{2} = \Setabs{p \in \Rat}{p^2 < 2\text{
अथवा }p < 0}$। निश्चय ही हमें जाँचना होगा कि यह सचमुच एक विच्छेद \emph{है},
पर यह जाँच \olref[check]{sec} में छोड़ते हैं।

पहले की तरह, कुछ वस्तुएँ परिभाषित करने के बाद उन पर मूल फलन और संबंध
परिभाषित करने हैं। एक सरल संबंध से आरंभ करें:
\begin{align*}
	\alpha \leq \beta \text{ तभी और केवल तभी जब }\alpha \subseteq \beta
\end{align*}
क्रम की यह परिभाषा हमें केंद्रीय परिणाम \emph{कहने} देती है: विच्छेदों के
समुच्चय में पूर्णता गुणधर्म है। पूरा खोलकर कहें तो कथन का रूप यह है। यदि $S$
विच्छेदों का कोई रिक्तेतर समुच्चय है और उसका कोई ऊपरी परिबंध है, तो $S$ का
एक लघुतम ऊपरी परिबंध है। अधिक विस्तार से: एक विच्छेद $\lambda$ है जो $S$ का
ऊपरी परिबंध है, अर्थात\ $(\forall \alpha \in S)\alpha \subseteq
\lambda$, और $\lambda$ ऐसा लघुतम विच्छेद है, अर्थात\ $(\forall \beta \in \Real)((\forall \alpha \in S)\alpha \subseteq \beta \lif \lambda \subseteq \beta)$। अब इस परिणाम का प्रमाण प्रस्तुत है:

\begin{thm}\ollabel{realcompleteness}
विच्छेदों के समुच्चय में पूर्णता गुणधर्म होता है।
\end{thm}

\begin{proof}
मान लें कि $S$ विच्छेदों का कोई रिक्तेतर समुच्चय है जिसका एक ऊपरी परिबंध है।
मान लें कि $\lambda
= \bigcup S$।
%\Setabs{p \in \Rat}{(\exists \alpha \in S)p \in \alpha}$$
पहले हमारा दावा है कि $\lambda$ एक विच्छेद है:
\begin{enumerate}
\item चूँकि $S$ का ऊपरी परिबंध है, इसलिए कम-से-कम एक विच्छेद $S$ में है, अतः
$\emptyset \neq \lambda$। चूँकि $S$ विच्छेदों का समुच्चय है, इसलिए $\lambda
\subseteq \Rat$। चूँकि $S$ का ऊपरी परिबंध है, कोई $p \in \Rat$ प्रत्येक
विच्छेद $\alpha \in S$ से अनुपस्थित है। अतः $p\notin \lambda$, और इसलिए
$\lambda \subsetneq \Rat$।
\item मान लें कि $p < q \in \lambda$। तब कोई $\alpha \in S$ ऐसा है कि
$q \in \alpha$। चूँकि $\alpha$ एक विच्छेद है, इसलिए $p \in \alpha$। अतः
$p \in \lambda$।
\item मान लें कि $p \in \lambda$। तब कोई $\alpha \in S$ ऐसा है कि
$p \in \alpha$। चूँकि $\alpha$ एक विच्छेद है, इसलिए कोई $q \in
\alpha$ ऐसा है कि $p < q$। अतः $q \in \lambda$।
\end{enumerate}
इससे दावा सिद्ध होता है। इसके अतिरिक्त स्पष्टतः $(\forall \alpha \in S)\alpha
\subseteq \bigcup S = \lambda$, अर्थात\ $\lambda$, $S$ का ऊपरी परिबंध है। अब मान लें कि $\beta \in \mathbb{R}$ भी एक ऊपरी परिबंध है, अर्थात\ $(\forall \alpha \in S)\alpha \subseteq \beta$। किसी भी $p \in \Rat$ के लिए, यदि $p \in \lambda$, तो कोई $\alpha \in S$ ऐसा है कि $p \in \alpha$, इसलिए $p \in \beta$। इसे व्यापक करने पर $\lambda \subseteq \beta$। अतः $\lambda$, $S$ का \emph{लघुतम} ऊपरी परिबंध है।
\end{proof}

इस प्रकार हमारे पास ऐसी बहुत-सी वस्तुएँ हैं जो पूर्णता गुणधर्म को संतुष्ट करती
हैं। इसे इस तरह भी कह सकते हैं: हमारे विच्छेदों में कोई ``रिक्त स्थान'' नहीं
है। (इसलिए परिमेय संख्याओं के बजाय वास्तविक संख्याओं के और ``विच्छेद'' लेने
से कोई रोचक नई वस्तु नहीं मिलेगी।)

अब वास्तविक संख्याओं पर कुछ संक्रियाएँ परिभाषित करनी हैं। पहले यह निर्धारित
करके कि $p_\Real =
\Setabs{q \in \Rat}{q < p}$, प्रत्येक $p \in \Rat$ को वास्तविक संख्याओं में
अंतःस्थापित करें। फिर परिभाषित करें:
\begin{align*}
	\alpha + \beta &= \Setabs{p + q}{p \in \alpha \land q \in \beta}\\
%	\alpha - \beta &\defis \Setabs{p - q}{p \in \alpha \land q \in \Rat \setminus \beta}\\
	\alpha \times \beta &= 
	\Setabs{p \times q}{0 \leq p \in \alpha \land 0 \leq q \in \beta} \cup 0^\mathbb{R} & \text{यदि }\alpha, \beta \geq 0_\Real
%	\alpha \div \beta &\defis
%	\Setabs{\nicefrac{p}{q}}{p \in \alpha \land q \in \Rat \setminus \beta}  & \text{if }\alpha \geq 0_\Real, \beta > 0_\Real
\end{align*}
गुणन के अन्य मामलों को सँभालने के लिए पहले रखें: %that $0_\Real \times \alpha = 0_\Real = \alpha \times 0_\Real$, and add:
\begin{align*}
	-\alpha &= \Setabs{p - q}{p < 0 \land q \notin \alpha}
\end{align*}
और फिर निर्धारित करें:
\begin{align*}
	\alpha \times \beta &\defis 
	\begin{cases}
		\mathord{-}\alpha \times \mathord{-}\beta &\text{यदि }\alpha < 0_\Real\text{ और }\beta < 0_\Real\\
		\mathord{-}(\mathord{-}\alpha \times \beta) &\text{यदि }\alpha < 0_\Real \text{ और }\beta > 0_\Real\\
		\mathord{-}(\alpha \times \mathord{-}\beta) &\text{यदि }\alpha > 0_\Real \text{ और }\beta < 0_\Real
	\end{cases}
\end{align*}
फिर जाँचना होगा कि इनमें से प्रत्येक परिभाषा का परिणाम हमेशा एक विच्छेद होता
है। अंत में एक सरल (पर लंबी) स्थापना करनी होगी कि इस प्रकार परिभाषित विच्छेद
ठीक वैसा ही व्यवहार करते हैं जैसा उन्हें करना चाहिए। ये सभी जाँचें
\olref[check]{sec} में छोड़ते हैं।
\end{document}
